% Part: first-order-logic
% Chapter: introduction
% Section: soundness-completeness

\documentclass[../../../include/open-logic-section]{subfiles}

\begin{document}

\olfileid{fol}{int}{scp}

\olsection{صحت مندی تے تمامیت}

اسیں پہلے درجے دی منطق لئی
!!{derivation} نظام وی متعارف
کراں گے۔ منطق داناں نے کئی
!!{derivation} نظام بنائے نیں،
پر اوہ سارے !!{sentence}s دے
درمیان اکّو !!{derivability}
تعلق تعریف کردے نیں۔ اسیں کہندے
آں کہ $\Gamma$ توں
~$!A$ \emph{!!{derive}s}
اے، $\Gamma \Proves !A$،
جے اک مخصوص، ٹھیک ٹھیک
تعریف کیتی قسم دا
!!a{derivation} موجود ہووے۔
!!^{derivation}s ہمیشہ
علامتاں دیاں محدود
بناوٹاں ہوندیاں نیں---
ہو سکدا اے !!{sentence}s
دی فہرست، یا کوئی ہور
گنجل دار ساخت۔ !!{derivation}
نظاماں دا مقصد ایہ طے
کرن دا وسیلہ دینا اے کہ
کی کوئی !!a{sentence}
کسے سیٹ~$\Gamma$ توں
منطقی طور تے نکلدا اے۔
ایس مقصد لئی ضروری اے
کہ $\Gamma \Entails !A$
اُدوں تے صرف اُدوں
ہووے جدوں $\Gamma
\Proves !A$ ہووے۔

جے $\Gamma \Proves !A$
ہووے پر $\Gamma \Entails !A$
نہ ہووے، تاں ساڈا
!!{derivation} نظام حد توں
ودھ طاقت ور ہووے گا:
اوہ بہت کجھ ثابت کر دے
گا۔ ایہ خاصیت کہ جے
$\Gamma
\Proves !A$ تاں
$\Gamma \Entails !A$
وی، \emph{صحت مندی}
کہلاندی اے، تے چنگے
!!{derivation} نظام لئی
ایہ گھٹ توں گھٹ شرط
اے۔ دوجے پاسے، جے
$\Gamma \Entails !A$
ہووے پر $\Gamma \Proves !A$
نہ ہووے، تاں ساڈا
!!{derivation} نظام حد توں
گھٹ طاقت ور ہووے گا:
اوہ کافی کچھ ثابت نہیں
کردا۔ ایہ خاصیت کہ جے
$\Gamma \Entails !A$
تاں $\Gamma \Proves !A$
وی، \emph{تمامیت}
کہلاندی اے۔ صحت مندی
دا ثبوت عموماً نسبتاً
آسان اے (!!{derivation}s
دی بناوٹ اُتے استقرا نال،
کیوں جے انہاں دی تعریف
استقرائی اے)۔ تمامیت دا
ثبوت اوکھا اے۔

صحت مندی تے تمامیت دے
کئی اہم نتیجے نیں۔ جے
!!{sentence}s دا کوئی
سیٹ~$\Gamma$ تضاد
!!{derive}s کرے (مثلاً
$!A \land \lnot !A$)،
تاں اوہ \emph{متضاد}
کہلاندا اے۔ متضاد
$\Gamma$s دے کوئی
نمونے نہیں ہو سکدے؛
اوہ پورا ہون جوگ
نہیں۔ تمامیت توں
اُلٹی گل نکلدی اے:
ہر اوہ $\Gamma$
جیہڑا متضاد نہیں---
یا جیویں اسیں آکھاں
گے، \emph{بے تضاد}---
اوہدا اک نمونہ ہوندا
اے۔ اصل وچ ایہ تمامیت
دے ہم ارز اے، تے
اسی روپ وچ اسیں
تمامیت ثابت کراں گے۔
ایہ ڈونگھا، شاید
حیران کن نتیجہ اے:
صرف ایس گل توں کہ
تسی $!A \land \lnot !A$
نوں $\Gamma$ توں
ثابت نہیں کر سکدے،
ایہ یقینی بن جاندا
اے کہ کوئی
!!a{structure} موجود
اے جیہڑی~$\Gamma$
دے بیان مطابق اے۔
ایس لئی تمامیت
ایس سوال دا جواب
دیندی اے: !!{sentence}s
دے کیہڑے سیٹاں دے
نمونے ہوندے نیں؟
جواب: اوہ سارے
تے صرف اوہی سیٹ
جیہڑے بے تضاد نیں۔

صحت مندی تے تمامیت
دے قضیاں دے دو
اہم نتیجے نیں: محدود
گواہی تے لووینہائم--سکولم
دا قضیہ۔ نمونیاں دے
نظریے وچ ایہ اہم
نتیجے نیں تے انہاں
نال کئی دلچسپ نتیجے
قائم کیتے جا سکدے
نیں۔ اسیں دو دا
ذکر پہلے کر چکے آں:
پہلے درجے دی منطق
ایہ بیان نہیں کر
سکدی کہ کسے
!!a{structure} دا
!!{domain} محدود
اے، یا اوہدا دائرہ
!!{nonenumerable}
اے۔

تاریخی طور تے ایہ
سب کجھ---پہلے درجے
دی منطق دی نحو تے
معنویات دی تعریف،
چنگے !!{derivation}
نظاماں دی تعریف،
انہاں دی صحت مندی
تے تمامیت دا ثبوت،
تے پہلے درجے دیاں
زباناں وچ کیہ بیان
ہو سکدا اے تے کیہ
نہیں، ایس دی وضاحت---
سمجھ کے درست کرن
وچ لمّا سماں لگیا۔
ہن سانوں ایہ کرنا
آؤندا اے، پر ساریاں
تفصیلاں وچوں لنگھنا
اج وی الجھن بھریا
تے تھکا دینے والا
ہو سکدا اے۔ پر
ایہ اہم وی اے،
کیوں جے پہلے درجے
دی منطق دی رسمی
زبان لئی ایتھے
بنائے طریقے منطق،
کمپیوٹر سائنس
تے لسانیات وچ
ہر پاسے ورتے
جاندے نیں۔ ایس لئی
تفصیلاں اُتے محنت
دا فائدہ لمبے
سمے وچ ملدا اے۔

\end{document}
